\section*{第 \ref{ch:g2:measuring-length} 章——测量长度}

\begin{solution}{exo:g2:measuring-length:1}
靠得近：背靠背站好。离得远：用\omterm{def:g2:measuring-length:units}{厘米}量出自己的身高，
然后把数字寄去——数字可以装进信里，
背靠背的比较却不行。
\end{solution}

\begin{solution}{exo:g2:measuring-length:2}
两人都没有数错。玛拉的步子较小，
所以同一间教室里能数出更多步。不同的\omterm{def:g2:measuring-length:measure}{单位}
得到的数字不同，但量的是同一段长度。
\end{solution}

\begin{solution}{exo:g2:measuring-length:3}
\omterm{def:g2:measuring-length:units}{厘米}对每个人都一样——它不属于任何人的
身体。用\omterm{def:g2:measuring-length:units}{厘米}得到的\omterm{def:g2:measuring-length:measure}{测量}结果，任何人
用任何一把尺都能核对。
\end{solution}

\begin{solution}{exo:g2:measuring-length:4}
蚂蚁：\unit{cm}（更小！）；门：\unit{m}；
操场：\unit{m}；这本书的宽度：\unit{cm}。
\end{solution}

\begin{solution}{exo:g2:measuring-length:5}
不对——尺子的边缘与零刻度线之间通常有一小段空白，
所以铅笔比 $13$ \omterm{def:g2:measuring-length:units}{厘米}短一点。
蒂姆忘了把\emph{零刻度线}而不是尺子边缘对准铅笔的
末端。
\end{solution}

\begin{solution}{exo:g2:measuring-length:6}
答案不唯一；常见的有：手掌宽约 \qty{7}{cm}，勺子长约
\qty{15}{cm}，杯子高约 \qty{9}{cm}。关键在于每个数字后面
都要带上 \unit{cm}。
\end{solution}

\begin{solution}{exo:g2:measuring-length:7}
$100 + 100 = 200$：门高 \qty{200}{cm}。
\end{solution}

\begin{solution}{exo:g2:measuring-length:8}
答案不唯一；大多数床的长度接近 \qty{2}{m}。
重点在于把估计值与\omterm{def:g2:measuring-length:measure}{测量}得到的数字进行比较。
\end{solution}

\begin{solution}{exo:g2:measuring-length:9}
\qty{45}{cm} 的丝带更长：$45 - 38 = 7$，所以长 \qty{7}{cm}。
不用数字，把两条丝带并排放在一起，也能分辨出哪条更长，
——但“\qty{45}{cm}”可以写进信里，而“并排放”
却不行。
\end{solution}

\begin{solution}{exo:g2:measuring-length:10}
没有别人的步长跟爷爷一样：一个小孩在同一个花园里走，
可能数出 $60$ 小步，而访客若没有爷爷在场，就无法核对“$30$
步”。用米来量——对每个人都一样——花园的长度
可以写下来、寄出去，
并由任何人核对。
\end{solution}
